\documentclass[10pt]{amsart}
\usepackage[a4paper,margin=23mm]{geometry}
\usepackage{amsmath,amssymb,amsthm,xfrac,hyperref,xspace,graphicx}
\usepackage[english]{babel}
\usepackage[scr=rsfs,cal=euler]{mathalfa}
\newcommand{\resp}{resp.}
\newcommand{\cf}{cf.}
\newcommand{\ie}{i.e.}
\newcommand{\eg}{e.g.}
\newcommand{\loccit}{loc.\ cit.}
\newcommand{\skpt}{\par\smallskip\noindent}
\let\Subsection\subsection
\let\Subsubsection\subsubsection
\emergencystretch=3em
\pagestyle{plain}
\hyphenation{irre-du-ci-ble dyna-mi-cal horo-grad-ing restric-tion deduce action expand-ing belong}

\newcommand{\Changel}{\mathcode`l="7160}
\newcommand{\Changelback}{\mathcode`l="716C}
\newcommand{\NoChangel}[1]{%
\expandafter\let\csname old\string#1\endcsname=#1
\let#1=\relax
\newcommand{#1}{\mathop{\mathcode`l="716C\csname old\string#1\endcsname\mathcode`l="7160}\displaylimits}%
}

\NoChangel{\log}
\NoChangel{\ln}
\NoChangel{\lim}
\NoChangel{\limsup}
\NoChangel{\liminf}
\NoChangel{\varlimsup}
\NoChangel{\varliminf}
\NoChangel{\varinjlim}
\NoChangel{\varprojlim}

\let\epsilon\varepsilon
\def\mto{\mathchoice{\longmapsto}{\mapsto}{\mapsto}{\mapsto}}

\usepackage{stmaryrd}

\newtheorem{thm}{Theorem}[section]
\newtheorem{introthm}{Theorem}
\renewcommand*{\theintrothm}{\Alph{introthm}}
\newtheorem*{claim}{Claim}
\newtheorem{claimnum}{Claim}
\newtheorem{lem}[thm]{Lemma}
\newtheorem{prop}[thm]{Proposition}
\newtheorem{cor}[thm]{Corollary}

\theoremstyle{definition}
\newtheorem{dfn}[thm]{Definition}
\newtheorem{rem}[thm]{Remark}
\newtheorem{ex}[thm]{Example}

\DeclareMathOperator{\Aff}{\mathsf{Aff}}
\DeclareMathOperator{\Aut}{\mathsf{Aut}}
\DeclareMathOperator{\Bij}{\mathsf{Sym}}
\DeclareMathOperator{\Br}{\mathsf{Br}}
\DeclareMathOperator{\BS}{\mathsf{BS}}
\DeclareMathOperator{\Diff}{\mathsf{Diff}}
\DeclareMathOperator{\Fit}{\mathsf{Fit}}
\DeclareMathOperator{\Fix}{\mathsf{Fix}}\let\fix\Fix
\DeclareMathOperator{\fixphi}{\mathsf{Fix}^\varphi}
\DeclareMathOperator{\homeo}{\mathsf{Homeo}}
\DeclareMathOperator{\Homirr}{\mathsf{Hom}_{\mathrm{irr}}}
\DeclareMathOperator{\hor}{\mathfrak{h}}
\DeclareMathOperator{\GL}{\mathsf{GL}}
\DeclareMathOperator{\id}{\mathsf{id}}
\DeclareMathOperator{\Int}{\mathsf{Int}}
\DeclareMathOperator{\stab}{\mathsf{Stab}}
\DeclareMathOperator{\supp}{\mathsf{Supp}}
\DeclareMathOperator{\suppphi}{\mathsf{Supp}^\varphi}
\DeclareMathOperator{\vf}{\mathsf{vf}}
\DeclareMathOperator{\wrwr}{{\wr\wr}}

\newcommand{\R}{\mathbb{R}}
\newcommand{\Q}{\mathbb Q}
\newcommand{\N}{\mathbb N}
\newcommand{\Z}{\mathbb Z}
\newcommand{\Tbb}{\mathbb T}

\newcommand{\Lcal}{\mathcal{L}}

\newcommand{\treeorder}{\mathrel{\triangleleft}}
\newcommand{\treeorderop}{\mathrel{\triangleright}}
\newcommand{\treeordereq}{\mathrel{\trianglelefteq}}
\newcommand{\treeup}{\mathrel{\curlywedgeuparrow}}

\begin{document}
\title{Plante-like wreath-product actions and their finitely generated epimorphic lifts cannot be semi-conjugate to C1 actions on the closed interval}
\author{Joaquín Brum \and Nicolás Matte Bon \and Cristóbal Rivas \and Michele Triestino}
\maketitle
\setcounter{section}{1}
\section{Notation and preliminaries}
\subsection{Actions on the line}\label{sc.preliminary-actions}
Let us start with some notation and terminology. When a group $G$ acts on a set $X$, we~simply write $g\cdot x$ for the action, unless there is risk of confusion. If~$\varphi\colon G\to \Bij(X)$ is an action, we~write $\stab^\varphi(x)$ for the stabilizer of $x\in X$ in $G$ under the action $\varphi$, $\fixphi(g)$ for the subset of fixed points in $X$ of an element $g\in G$, and $\suppphi(g)=X\setminus \fixphi(g)$ for its support.

Let $X$ be a real interval with non-empty interior (that is, up to diffeomorphism, one of the spaces $\R, [0,1], [0, 1)$). The interior of $X$ will be denoted by $\Int(X)$. The group of orientation-preserving homeomorphisms of $X$ is denoted by $\homeo_0(X)$, and similarly we denote the group of orientation-preserving $C^r$ diffeomorphisms $\Diff_0^r(X)$, for $r\ge 1$. We~will say that a group action $\varphi\colon G\to \homeo_0(X)$ is \emph{irreducible} if it has no fixed point in $\Int(X)$. We~also recall the notion of semi-conjugacy for actions on intervals (see Kim, Koberda, and Mj \cite{KKMj}).
\begin{dfn}
Let $X$ and $Y$ be real intervals with non-empty interior, and let $\varphi\colon G\to \homeo_0(X)$ and $\psi\colon G\to \homeo_0(Y)$ be two irreducible actions of a group $G$. We~say that $\varphi$ and $\psi$ are \emph{semi-conjugate} if there exists a monotone map $h\colon \Int(X)\to \Int(Y)$ such that $h\circ\varphi(g)=\psi(g)\circ h$ for any $g\in G$. When $h$ is non-decreasing, we~say that $\varphi$ and $\psi$ are \emph{positively} semi-conjugate. When $h$ is an (orientation-preserving) homeomorphism, we~say that $\varphi$ and $\psi$ are \emph{(positively) conjugate}.
\end{dfn}

A well-known argument (see Navas \cite[Prop.\,2.1.12]{Navas-book}) shows that when the group $G$ is finitely generated, every irreducible group action $\varphi\colon G\to \homeo_0(\R)$ admits a minimal non-empty closed invariant subset (that we shall call for short a \emph{minimal invariant subset}), which is either $\R$, a perfect subset of empty interior, or~a closed (and thus discrete) orbit. From the perspective of semi-conjugacy, this implies that $\varphi$ is always semi-conjugate to an action $\psi\colon G\to \homeo_0(\R)$ which is either \emph{minimal} (that is, every orbit is dense), or~which takes values in the group~$(\Z, +)$ of integer translations, in~which case we will say that $\psi$ is \emph{cyclic}. Moreover, such a minimal or cyclic representative of each semi-conjugacy class is unique up to conjugacy. Thus when studying continuous actions of $G$ up to semi-conjugacy it is enough to restrict to minimal actions (the cyclic case being rather trivial), and we will often do so.

We will often build actions of groups on real intervals starting from actions on totally ordered sets. This strategy is quite general, but for the purpose of this note we will restrict our attention to actions on ordered sets $(\Omega,<)$ which are countable, have neither minimal nor maximal element, and which are \emph{densely ordered} in the sense that for every $\omega_1<\omega_2$ in $\Omega$ there is $\omega\in \Omega$ such that $\omega_1<\omega<\omega_2$. The good thing about these sets is that Cantor's ``back and forth argument'' ensures that there is an order-preserving bijection $t\colon (\Omega,<)\to (\Q,<)$ where rationals are equipped with the natural order (see for instance Clay and Rolfsen \cite[Th.\,2.22]{ClayRolfsen}).

\begin{dfn} \label{def dynamical realization}
Let $(\Omega,<)$ be a countable densely ordered set having neither maximal nor minimal element, $t\colon \Omega\to \Q$ an order-preserving bijection, and $\varphi\colon G\to \Aut(\Omega,<)$ a non-trivial order-preserving action. A \emph{dynamical realization} of $\varphi$ is an irreducible action $\bar\varphi\colon G\to \homeo_0(\R)$ with the property that
$\bar\varphi(g)(t(\omega))= t(\varphi(g)(\omega))$
for every $g\in G$ and every $\omega\in \Omega$.
\end{dfn}
A dynamical realization always exists, see \cite[\S 2.4]{ClayRolfsen} for details. Moreover, if~$\bar\varphi_1$ and $\bar\varphi_2$ are two dynamical realizations of $\varphi\colon G\to \Aut(\Omega,<)$, then they are positively conjugate. Indeed, if~$t_1$ and $t_2$ are the corresponding bijections from $\Omega$ to $\Q$, then $h=t_1\circ t_2^{-1}$ can be extended to the whole real line so that it conjugates $\bar\varphi_1$ to $\bar\varphi_2$. This is why we usually refer to \emph{the} dynamical realization of $\varphi$. The following result will be useful later.

\begin{lem} \label{l-dynamical-realisation}
Let $(\Omega,<)$ be a countable densely ordered set with neither maximal nor minimal element, and $\varphi\colon G\to \Aut(\Omega,<)$ an order-preserving action whose dynamical realization $\bar\varphi\colon G\to \homeo_0(\R)$ is minimal.
Let $\psi\colon G\to \homeo_0(\R)$ be an irreducible action, and assume that $\sigma\colon \Omega\to \R$ is a map which is non-decreasing and $G$-equivariant, in~the sense that
\[ \sigma (\varphi(g)(\omega))=\psi(g)(\sigma(\omega)) \quad \quad \text{for every }g\in G, \omega\in \Omega.
\]
Then, $\sigma$ is injective, and $\psi$ and $\bar{\varphi}$ are positively semi-conjugate. In~particular, when $\psi$ is minimal, it~is positively conjugate to $\bar{\varphi}$.
\end{lem}

\begin{proof}
Let $t\colon \Omega\to \Q$ be the bijection used to define the dynamical realization $\bar\varphi$. We~define a map $\tilde{\sigma}\colon \R \to \R$ by setting $\tilde{\sigma}(x)=\sup \{\sigma(\omega) : t(\omega) \le x\}$. Then $\tilde{\sigma}$ is non-decreasing and equivariant with respect to the $\bar{\varphi}$-action on the source and the $\psi$-action on the target, showing that $\psi$ is positively semi-conjugate to $\bar{\varphi}$. We~claim that $\tilde{\sigma}$ is injective. To see this, consider the subset
\[S=\{x\in\R:\tilde\sigma\text{ is constant on a neighborhood of }x\}.
\]
It is direct to see that $S$ is open and $\bar\varphi$-invariant, and therefore empty by minimality of $\bar{\varphi}$. To conclude the proof of the claim, notice that if $\tilde{\sigma}(x)\neq\tilde\sigma(y)$ for $x<y$, then monotonicity of $\tilde\sigma$ would imply $(x,y)\subseteq S$, which is a contradiction. In~particular, since $t$ is injective, so is $\sigma=\tilde{\sigma}\circ t$.

Assume now that $\psi$ is minimal. In this case, the image of $\tilde\sigma$ must be dense, or~otherwise its closure would be a proper closed $\psi$-invariant subset. Therefore, since $\tilde{\sigma}$ is monotone, injective, and has dense image, it~must be a homeomorphism. This shows that $\bar{\varphi}$ and $\psi$ are positively conjugate.
\end{proof}
\setcounter{subsection}{2}
\subsection{The group $\Z\wr\Z$ as an illustrative example} \label{sec ejemplo lamp lighter}
Throughout this subsection we set $\Z \wr \Z=L\rtimes \Z$, where $L=\bigoplus_\Z \Z$ is the group of finitely supported configurations (functions) $f\colon \Z\to \Z$, and $\Z$ acts on $L$ by shifting the indices. We~fix the generating set $\{g,h_0\}$ of $\Z\wr\Z$, where $g$ generates the factor~$\Z$, and $h_0\in L$ is the configuration given by $h_0(0)=1$ and $h_0(k)=0$ for $k\neq 0$, and denote by $h_n=g^nh_0g^{-n}$ ($n\in \Z$) the elements given by $h_n(n)=1$ and $h_n(k)=0$ for $k\neq n$.

The group $\Z\wr\Z$ is solvable (more precisely, metabelian) and it is clear that $L$ agrees with $\Fit(\Z\wr\Z)$, the Fitting subgroup of $\Z\wr\Z$. Moreover, $\Z\wr\Z$ admits the infinite presentation
\begin{equation}\Z\wr\Z\cong\langle g, h_0 \mid h_n=g^nh_0g^{-n}, \, [h_n, h_m]= \id \quad (n, m\in \Z)\rangle. \label{e-presentation-ZwrZ}\end{equation}

\subsubsection{Affine actions}
Another way of producing actions of $\Z\wr\Z$ on the line is by considering affine maps
\[\varphi(g) \colon x\mto \lambda x+\beta , \quad \varphi(h_0)\colon x\mto x+ \alpha \quad \quad (\lambda>0\text{ and } \alpha, \beta \in\R).
\]
The reader can check that for any choice of $\lambda,\alpha,\beta$ as above, the images of the generators satisfy the relations from \eqref{e-presentation-ZwrZ}, and hence the action $\varphi\colon\Z\wr\Z\to \homeo_0(\R)$ is well defined. These affine actions can be readily classified up to conjugacy.

\begin{prop} \label{l-wr-affine} Let $\{g,h_0\}$ be the generating set for $\Z\wr\Z$ as above. Then every irreducible affine action $\varphi\colon \Z \wr \Z\to \Aff(\R)$ is positively conjugate by an affine map to an action in one of the following families:
\begin{itemize}
\item actions by translations obtained by setting
\[\varphi(h_0)\colon x\mto x+\alpha, \quad \varphi(g)\colon x\mto x +\beta \quad \quad (\alpha^2+\beta^2=1);
\]
\item non-abelian affine actions obtained by setting
\[\varphi(h_0)\colon x\mto x\pm 1, \quad \varphi(g) \colon x\mto \lambda x\quad \quad (\lambda>0).
\]
\end{itemize}
\end{prop}
\begin{proof}
Let $\varphi$ be an irreducible affine action of $\Z\wr\Z$. First we show that $\varphi(h_0)$ is a translation. Assume by contradiction that $\varphi(h_0)$ is a homothety, with fixed point $p$. Then $p$ cannot be fixed by $\varphi(g)$ or it would be a global fixed point. Thus $\varphi(h_1)$ is a homothety whose fixed point is $\varphi(g)(p)\neq p$. This contradicts that $\varphi(h_0)$ and $\varphi(h_1)$ commute. Thus $\varphi(h_0)$ is a translation (possibly trivial). If~$\varphi(g)$ is also a translation, up to conjugation by an affine map we are in the first case of the classification. If~$\varphi(g)$ is a homothety, then $\varphi(h_0)$ cannot be trivial, and upon conjugating by an affine map we can assume that $\varphi(h)$ is a translation by $1$ and that $\varphi(g)$ fixes 0, which yields the second case of the classification.
\end{proof}

\begin{rem}
In the previous classification, note that non-abelian affine actions are faithful if and only if $\lambda$ is a transcendental number (since the translations $\varphi(h_n)(x)=x+\lambda^n$ generate a free abelian group precisely when $\lambda $ is transcendental). Actions by translations obviously descend to actions of the abelianization of $\Z \wr \Z$, which is isomorphic to $\Z^2$. The abelianization acts faithfully provided $\alpha$ and $\beta$ are rationally independent.
\end{rem}
\subsubsection{Plante-like actions}
\label{ssc.plante} Finally, we~describe a procedure for producing actions of $\Z\wr\Z$ on the line that are not semi-conjugate to any (real) affine action. Up to positive conjugacy, there will be four such actions (two, if~we also allow negative conjugacies) which are what we call \emph{Plante-like actions} of $\Z\wr\Z$. One of these actions was first considered by Plante in \cite{Plante}, who constructed it by explicitly defining two homeomorphisms of the line. Following the discussion in our previous work \cite[Ex.\,8.1.8]{BMRT}, we~shall take a different point of view, and obtain such actions as the dynamical realization of an affine action of $\Z\wr\Z$ on a countable ordered set (this point of view can be easily generalized to other wreath products of groups, as explained below).

Recall that we write $\Z \wr \Z=L\rtimes \Z=\langle g,h_0\rangle$, where $L=\bigoplus_\Z \Z$ and $\{ g, h_0\}$ from \eqref{e-presentation-ZwrZ} is our preferred generating set, and that we set $h_n=g^n h_0 g^{-n}$.

One can identify the subgroup $L$ with the ring of Laurent polynomials $\Z[X,X^{-1}]$, choosing the identification so that $h_n\in L$ represents the polynomial $ X^n$ (in~particular, $h_0$ represents the constant polynomial $1$). With this point of view, the group $\Z\wr\Z$ admits a natural action on $\Z[X,X^{-1}]$, where $L\cong \Z[X,X^{-1}]$ acts additively and the generator $g$ acts by multiplication by $X$. This action of $\Z \wr \Z$ on $\Z[X,X^{-1}]$ preserves four natural orders $\prec_{\max}^\pm$ and $\prec_{\min}^\pm$ of lexicographic type, obtained by looking at the sign of the coefficient of the monomial of maximal or minimal degree, respectively. More precisely, given a non-zero Laurent polynomial
$P(X)$, we~declare $0\prec_{\max}^+ P(X)$ if the coefficient of the highest power of $X$ in $P(X)$ is a non-negative integer, and $0\prec_{\max}^- P(X)$ if the coefficient of the highest power of $X$ in $P(X)$ is a non-positive integer. Similarly, we~declare $0\prec_{\min}^+ P(X)$ (respectively, $0\prec_{\min}^- P(X)$) if the coefficient of the lowest power of $X$ in $P(X)$ is a non-negative integer (respectively, non-positive integer). It~is routine to check that these four orders are invariant under the action of $\Z \wr \Z$ defined above. By~considering the dynamical realizations of the above actions, we~obtain, up to positive conjugacy, four actions on the line which are what we call the \emph{Plante-like actions} of $\Z \wr\Z$.

Let us take a closer look at Plante-like actions. Let $\prec$ be one of the orders $\prec_{\max}^\pm$, $\prec_{\min}^\pm$ on $\Z[X, X^{-1}]$, and $\varphi\colon G\to \homeo_0(\R)$ the corresponding Plante-like action. On the one hand, by~definition of the lexicographic order $\prec$, for every $h\in L$ we have that every $h$-orbit in $(\Z[X,X^{-1}], \prec)$ is bounded. For example, in~the case of $\prec_{\max}^+$, the $h$-orbit of any $Q\in \Z[X, X^{-1}]$ is contained in the interval $(-X^m, X^m)$ whenever $m$ is any integer strictly larger than the degree of $Q$ and of the polynomial representing~$h$. However, the global action of the group $L$ on $\Z[X, X^{-1}]$ is transitive. This implies that for the Plante-like action $\varphi$, the group $L$ acts without fixed points (in~particular, $\varphi$ is irreducible), while every $h\in L$ has infinitely many fixed points. This immediately prevents the action $\varphi$ to be semi-conjugate to any affine action (cf.\ Proposition \ref{l-wr-affine}).

Consider now the action of the generator $g$ on $(\Z[X, X^{-1}], \prec)$. Since $g$ acts by multiplication by $X$, it~has exactly one fixed point (the polynomial $0$). Suppose, for definiteness, that $\prec$ is one of the orders $\prec_{\max}^\pm$. Then for any Laurent polynomial $P(X)\succ 0$, we~have $g\cdot P(X)=XP(X)\succ P(X)$, and moreover $g^n\cdot P(X)=X^nP(X)$ is arbitrarily large as $n\to +\infty$, and arbitrarily close to $0$ as $n\to -\infty$ (this easily follows from the definition of the lexicographic order, since applying $g^n$ to $P(X)$ shifts the largest power of $X$ by $n$, without changing the corresponding coefficient). Similarly for $P(X)\prec 0$, we~have that $g^n\cdot P(X)$ gets arbitrarily small as $n\to +\infty$, and arbitrarily close to $0$ as $n\to -\infty$. We~deduce that in the Plante-like action, $\varphi(g)$ is an expanding homothety (in the sense of Definition \ref{d-expanding-homothety}). Similarly one shows that $\varphi(g)$ is a contracting homothety in the case where $\prec$ is $\prec_{\min}^\pm$. We~can (and will) assume that the unique fixed point of $\varphi(g)$ is the origin of the real line. In Figure \href{https://doi.org/10.5802/jep.284}{1 of the source article} the Plante-like action corresponding to $\prec_{\max}^+$ is depicted.

\setcounter{thm}{7}
After these considerations, we~can prove the following.

\begin{prop} A Plante-like action of $\Z\wr\Z$ is minimal.
\label{prop Plante-like minimal}
\end{prop}

\begin{proof} Let $x\in \R$ be a point and $I\subset \R$ an open interval. We~need to check that the orbit of $x$ meets the interval $I$.
Since $I$ is open, it~follows that it contains rational points and hence points in the image of the order-preserving bijection $t\colon \Z[X,X^{-1}]\to \Q$ from which the dynamical realization is built (see Definition \ref{def dynamical realization} and the discussion before it). In~particular, since $L$ acts transitively on the image of $t$, there is $h\in L$ such that $\varphi(h)(I)$ contains $0\in \R$. Now, recall that we are assuming that the origin $0$ is the unique fixed point of $\varphi(g)$, which moreover acts as a homothety (expanding or contracting, depending on the order considered on $\Z[X,X^{-1}]$). It~follows that there is $n\in \Z$ such that $\varphi(g^nh)(I)$ contains the point $x$, and hence $\varphi(h^{-1}g^{-n})(x)$ belongs to $I$, as desired.
\end{proof}

\section{Lamination notation}\label{sec generalities}
\subsection{Laminations}
We say that two open intervals $I,J\subset \R$ \emph{do not cross} if they are either nested or disjoint, that is either $I\subset J$, or~$J\subset I$, or~$I\cap J=\varnothing$. A \emph{prelamination} of the real line is a collection $\Lcal$ of (non-empty) bounded open intervals, that pairwise do not cross. We~shall identify the set of all bounded open intervals in $\R$ with the space $\R^{(2)}:=\{(x, y)\in \R^2: x<y\}$, endowed with the topology induced from $\R^2$ (\ie convergence of endpoints). A \emph{lamination} is a prelamination which is moreover a closed subset of $\R^{(2)}$. A lamination is \emph{covering} if it defines a cover of $\R$; equivalently, if~it contains an increasing exhaustion of $\R$. Elements of a (pre)laminations will be often called \emph{leaves}.

Note that a prelamination $\Lcal$ is naturally a partially ordered set (\emph{poset}, for short) with respect to the inclusion order $\subseteq$. Given $l_0, l_1\in \Lcal$, we~will say that $l_1$ is the successor of $l_0$ if it is the smallest element of $(\Lcal, \subseteq)$ such that $l_0\subsetneq l_1$.

\begin{rem}
If an action $\varphi \colon G \to \homeo_0(\R)$ has an invariant prelamination $\Lcal$, then the closure $\overline{\mathcal L}$ in $\R^{(2)}$ is an invariant lamination. It~is convenient to visualize a prelamination of the line as a collection of pairwise disjoint semi-circles in the upper half-plane (\ie geodesic lines in the Poincaré metric), whose endpoints define the leaves of the prelamination.
\end{rem}

\begin{dfn}\label{def.lamination}
An irreducible action $\varphi\colon G\to\homeo_0(\R)$ is \emph{laminar} if it preserves a covering prelamination $\Lcal$. We~also say that the laminar action $\varphi$ is \emph{focal} if, in~addition, there exists $l\in\Lcal$ and a sequence $(g_n)\subset G$ such that $(\varphi(g_n)(l))$ is an increasing exhaustion of the line. When we want to emphasize the role of the prelamination, we~will say that $\varphi$ is laminar (and focal) \emph{with respect to $\Lcal$}.
\end{dfn}

\begin{rem}
The property of being a laminar (focal) action is invariant under semi-conjugacy, and every focal laminar action is semi-conjugate to a \emph{minimal} focal laminar action \cite[Prop.\,8.1.14]{BMRT}.
\end{rem}
\begin{rem}\label{r-minimal-implies-focal}
Conversely, assume that $\varphi\colon G\to\homeo_0(\R)$ is a {minimal} action. If~$\Lcal$ is any non-empty $\varphi$-invariant prelamination, then $\Lcal$ is automatically covering, and $\varphi$ is focal with respect to $\Lcal$ \cite[Prop.\,8.1.15]{BMRT}.
\end{rem}

Laminar actions satisfy some general constraints. In~particular, group elements can be classified according to two types of dynamics. For this, we~recall the following terminology.

\begin{dfn} \label{d-expanding-homothety}
A homeomorphism $h\in\homeo_0(\R)$ is an \emph{expanding pseudo-homothety} if there exists a compact subset $K\subset \R$ such that $h(U)\supset \overline{U}$ for every open subset $U\supset K$.
When $h^{-1}$ is an expanding pseudo-homothety, we~say that~$h$ is a \emph{contracting pseudo-homothety}. In the case where $K$ can be taken to be a point, we~simply say $h$ is an (expanding or contracting) \emph{homothety}. We~also say that $h\in \homeo_0(\R)$ is \emph{totally bounded} if $\Fix(h)$ accumulates on both $+\infty$ and $-\infty$.
\end{dfn}
We now recall the notion of horograding from \cite{BMRT}.
Let $\Lcal$ be a prelamination of the real line. A \emph{horograding} of $\Lcal$ is a monotone map $\hor \colon (\Lcal, \subseteq) \to (\R, \leq)$, namely a map such that $\hor(l_1)\leq \hor(l_2)$ whenever $l_1\subseteq l_2$ (in this case we say that $\hor$ is a \emph{positive} horograding) or $\hor(l_1)\geq \hor(l_2)$ whenever $l_1\subseteq l_2$ (in this case we say that $\hor$ is \emph{negative}). In the presence of a group action, this leads to the following notion.

\begin{dfn}
Let $\varphi\colon G\to \homeo_0(\R)$ be a laminar action, and $j\colon G\to \homeo_0(\R)$ an irreducible action. A (positive or negative) \emph{horograding of $\varphi$ by $j$} is a pair $(\Lcal, \hor)$ consisting of a $\varphi$-invariant covering lamination, and a (positive or negative, accordingly) horograding $\hor\colon (\Lcal, \subseteq) \to \R$ such that $\hor(\varphi(g)(l))=j(g)(\hor(l))$, for every $g\in G$ and $l\in \Lcal$.
\end{dfn}

\begin{rem}[{\cite[Lem.\,8.2.7]{BMRT}}] \label{r-horograding-proper}
Suppose that $(\Lcal, \hor)$ is a positive horograding of a laminar action $\varphi\colon G\to \homeo_0(\R)$ by an irreducible action $j\colon G\to \homeo_0(\R)$. Then for every increasing sequence $(l_n)\subset \Lcal$ that exhausts $\R$, we~must have $\hor(l_n)\to +\infty$. Indeed, since $j$ is irreducible and $\hor(\Lcal)$ is $j$-invariant, for every $M>0$ there exists $l\in \Lcal$ such that $\hor(l)\ge M$, and since $l_n$ eventually contains any given $l$, the conclusion follows.
\end{rem}

The existence of a horograding of a laminar action $\varphi$ by an action $j$ implies that the action $j$ retains information on the large-scale behavior of $\varphi$. In~particular, the type of each element with respect to $\varphi$ is determined by $j$ as follows.

\begin{prop}[{\cite[Prop.\,8.2.10]{BMRT}}]\label{p-classification-elements} Let $\varphi\colon G\to\homeo_0(\R)$ be a minimal laminar action positively horograded by $j\colon G\to\homeo_0(\R)$. For any element $g\in G$, we~have the following.
\begin{enumerate}
\item \label{i-pseudo-homothety} If $\Fix(j(g))$ does not accumulate on $+\infty$, then $\varphi(g)$ is a pseudo-homothety, which is expanding if $j(g)(x)>x$ for any sufficiently large $x$, and contracting otherwise.
\item Otherwise, $\varphi(g)$ is totally bounded.
\end{enumerate}
Moreover, in~the former situation, if~$\Fix(j(g))=\emptyset$ then $\varphi(g)$ is a homothety.
\end{prop}

\setcounter{thm}{9}
\Subsection{The example of Plante-like actions} \label{s-plante-laminations}
\subsubsection{Plante-like actions of $\Z \wr \Z$}
Here we explain why the Plante-like actions of $\Z \wr \Z$ (see Section~\ref{ssc.plante}) are laminar and horograded by a cyclic action.

\begin{prop}\label{prop Plante is laminar}
A Plante-like action $\varphi\colon \Z \wr \Z \to \homeo_0(\R)$ is laminar and horograded by a cyclic action.
\end{prop}

\begin{proof}
We resume notation from Section~\ref{sec ejemplo lamp lighter}. In~particular, recall that we denote by $\{g, h_0\}$ the standard generating pair of $\Z \wr \Z$, satisfying the presentation \eqref{e-presentation-ZwrZ}, and set $h_n:=g^n h_0 g^{-n}$.
For $n\in \Z$, we~let $\Lcal_n$ be the set of connected components of $\suppphi(h_n)$, and $\Lcal=\bigcup_n \Lcal_n$. For $l\in \Lcal_n$, write $\hor(l):=n$. We~want to prove that $\Lcal$ is a $\varphi$-invariant lamination (if so, since $\varphi$ is minimal by Proposition \ref{prop Plante-like minimal}, Remark~\ref{r-minimal-implies-focal} guarantees that it is laminar), and that $(\Lcal, \hor)$ is a horograding of $\varphi$ by the cyclic action $j\colon \Z \wr \Z\to \homeo_0(\R)$, given by $j(g):x\mto x+1$ and $j(h_0)=\id$.

We keep denoting by $L=\bigoplus_\Z \Z\cong \Z[X, X^{-1}]$ the group of lamps.
Recall from Section~\ref{ssc.plante} that there are four Plante-like actions of $\Z \wr \Z$, associated with four lexicographic orders on $\Z[X, X^{-1}]$. We~only discuss the case of the order $\prec_{\max}^+$ (the other cases are analogous), and start with the following observation (that we isolate for further reference).

\begin{claimnum} \label{l-plante-fixed-points}
For every $n\in \Z$, we~have $\suppphi(h_n)\subsetneq \suppphi(h_{n+1})$. More precisely, every connected component of $\suppphi(h_n)$ is compactly contained in $\suppphi(h_{n+1})$.
\end{claimnum}
\setcounter{claimnum}{0}
\begin{proof}[Proof of the claim]
Since each $h_n$ acts without fixed points on $\Z[X,X^{-1}]$, the open subset $\suppphi(h_n)$ contains $\Q$, which is the image of the order-preserving bijection $t\colon \Z[X,X^{-1}]\to \Q$ used to define the dynamical realization $\varphi$. More precisely, if~for every $P\in \Z[X, X^{-1}]$ we let $I(n, P)$ denote the connected component of $\suppphi(h_n)$ containing $t(P)$, we~have that every connected component of $\suppphi(h_n)$ is of the form $I(n, P)$, by~density of $\Q$. Note that each $I(n, P)$ is equal to the convex hull of $\{h_n^k\cdot t(P): k\in \Z\}$.

Recall that $h_n$ corresponds to the polynomial $X^n$ under the identification $L\cong \Z[X, X^{-1}]$. Using this identification, for $f_1, f_2\in L$, we~have that $f_1\cdot t(P)< f_2\cdot t(P)$ if and only if $f_1+P\preceq_{\max}^+ f_2+P$, which in turn is equivalent to $f_1\preceq_{\max}^+ f_2$. Now, by~the definition of the order $\preceq_{\max}^+$, we~have that
\[-X^{n+1} \preceq_{\max}^+ kX^n \preceq_{\max}^+ X^{n+1}
\]
for every $k\in \Z$. We~conclude that for every $P\in \Z[X, X^{-1}]$, we~have
\[h_{n+1}^{-1}\cdot t(P)< h_n^k\cdot t(P)< h_{n+1}\cdot t(P)
\]
for every $k\in \Z$, and thus ${I(n, P)}$ is compactly contained in $I(n+1, P)$.
\end{proof}

After Claim \ref{l-plante-fixed-points}, we~have that the collection $\Lcal=\bigcup_n \Lcal_n$ defined above is a prelamination. Next, we~check $\varphi$-invariance: by commutativity, the image of the generator $\varphi(h_0)$ preserves each support $\suppphi(h_n)$ by commutativity, so each $\Lcal_n$ is $\varphi(h_0)$ invariant; for the other generator, we~have
\[g\cdot \suppphi(h_n)=\suppphi(gh_ng^{-1})=\suppphi(h_{n+1}),
\]
so $g\cdot \Lcal_n=\Lcal_{n+1}$. This also gives that the map $\hor$ intertwines the $\varphi$-action on $\Lcal$ with the cyclic action $j$ defined above, defining a positive horograding.

Finally, let us check that $\Lcal$ is a discrete (hence closed) subset of $\R^{(2)}$. Indeed, suppose that $(l_n)\subset \Lcal$ is a sequence converging to some interval $I\in \R^{(2)}$. Up to discard finitely many terms, this gives that the intervals $l_n$ are all related by inclusion, and upon extracting a subsequence, we~can assume that the sequence $(l_n)$ is monotone, and thus also $(\hor(l_n))$ is. Now, if~the sequence $\hor(l_n)$ is bounded, then Claim \ref{l-plante-fixed-points} implies that $(l_n)$ is eventually constant and equal to $I$, in~particular $I\in \Lcal$. If~$\hor(l_n)$ increases to $+\infty$, then $I=\bigcup l_n$ is a connected component of $\bigcup_{m\in \Z}\suppphi(h_m)$; however, as the latter is a $\varphi$-invariant open set, minimality of $\varphi$ (Proposition \ref{prop Plante-like minimal}) gives $I=\R$, contradicting that $I$ is a bounded open interval. Finally, if~$\hor(l_n)\to -\infty$, we~have that~$I$ is a connected component of the interior of $\bigcap_{m\in \Z}\suppphi(h_m)$; using minimality again, we~conclude that $I=\varnothing$. Thus $\Lcal$ is a discrete subset of $\R^{(2)}$.
\end{proof}
\setcounter{section}{6}
\section{\texorpdfstring{$C^1$}{C1} actions on intervals} \label{sc.C1}
\begin{prop}\label{p-plante-C1}
Let $\rho\colon H\to \Z\wr \Z$ be an epimorphism from a finitely generated group $H$, and let $\psi \colon H\to \homeo_0(\R)$ be the action obtained by postcomposing $\rho$ with a Plante-like action. Let $\varphi\colon H\to \homeo_0([0, 1])$ be an action semi-conjugate to $\psi$. Then, $\varphi$ cannot be of class $C^1$ on $[0,1]$.
\end{prop}
Before the proof, let us recall from Sections~\ref{ssc.plante} and \ref{s-plante-laminations} some general properties of Plante-like actions, that will be used in the proof.
As before, we~denote by~$g$ and~$h_0$ the standard generators of $\Z \wr \Z$, and set $h_n=g^n h_0g^{-n}$. For definiteness, let $\eta\colon \Z\wr\Z\to \homeo_0(\R)$ be the Plante-like action associated with the lexicographic order $\prec_{\max}^{+}$. Recall that we choose the order-preserving bijection $t\colon \Z[X, X^{-1}]\to \Q$ used to construct the Plante-like action in such a way that the polynomial 0 is sent to the origin $0\in \R$, and this implies that $\eta(g)$ is an expanding homothety with $0$ as fixed point. In contrast, for every $n\in \Z$, the element $h_n$ satisfies $\eta(h_n)(x)\ge x$ for every $x\in \R$, with strict inequality for $x=0$ (since $h_n$ corresponds to the polynomial $X^n\in \Z[X, X^{-1}]$, which is positive for the lexicographic order $\prec_{\max}^+$).

For $n\in \Z$, let $I_n$ be the connected component of $\supp^\eta(h_n)$ containing $0$. Recall from the proof of Proposition \ref{prop Plante is laminar} (in~particular, see Claim \ref{l-plante-fixed-points}) that each $I_n$ is a bounded interval belonging to an invariant lamination, and that $\overline{I}_n\subseteq I_{n+1}$. Note also that $\eta(g)(I_n)=I_{n+1}$ (since $gh_ng^{-1}=h_{n+1}$). We~will use the following property of the Plante-like action $\eta$.
\begin{lem} \label{l-disjoint-intervals}
With notation as above, for $n\in \Z$, consider the subgroup $L_n:=\langle h_j\, (j\geq n)\rangle$ of $\Z\wr \Z$.
If $f_1, f_2\in L_{-n}$ are distinct elements, then $\eta(f_1)(I_{-n-1})$ and $\eta(f_2)(I_{-n-1})$ are disjoint.
\end{lem}

\begin{proof}
It is enough to show that $\eta(f)(I_{-n-1})$ and$I_{-n-1}$ are disjoint whenever $f\in L_{-n}$ is non-trivial. Such an $f$ can be uniquely written as $f=h_{m_1}^{k_1}\cdots k_{m_r}^{k_r}$, with $m_1>\cdots >m_r\ge -n$ and $k_1,\ldots, k_r\neq 0$. Claim \ref{l-plante-fixed-points} in Proposition \ref{prop Plante is laminar} gives the chain of inclusions $I_{-n-1}\subset I_{m_1}\subset \cdots \subset I_{m_r}$. Since each $I_{m_i}$ is $\eta(h_{m_i})$-invariant, this implies
\[\eta(h_{m_2}^{k_2}\cdots h_{m_r}^{k_r})(I_{-n-1})\subset I_{m_2}\subseteq I_{m_1-1}.
\]
Now note that $\eta(h_{m_1})^{k_1}$ moves both endpoints of $I_{m_1-1}$, and in the same direction (depending on the sign of $k_1$); since $\eta(h_{m_1}^{k_1})(I_{m_1-1})$ and $I_{m_1-1}$ cannot cross, they must be disjoint. This gives the desired conclusion.
\end{proof}

Lemma \ref{l-disjoint-intervals} allows to produce plenty of disjoint images of $I_0$ using controlled group elements. On the other hand, if~the action $\eta$ were semi-conjugate to a $C^1$ action on $[0, 1]$, the length of these intervals can be bounded below using derivatives. This is the strategy to obtain a contradiction in the proof below.

\begin{proof}[Proof of Proposition \ref{p-plante-C1}]
As above we assume that $\eta\colon \Z\wr\Z\to \homeo_0(\R)$ is the Plante-like action associated with the lexicographic order $\prec_{\max}^{+}$, and that the action~$\psi$ in the statement is $\psi=\eta\circ \rho$ (the case of the other Plante-like actions can be treated similarly). We~choose and fix preimages $\widetilde{g}, \widetilde{h}_0\in H$ of the generators $g, h_0$ under the epimorphism $\rho$, and set $\widetilde{h}_n:=\widetilde{h}_0$. Let $\varphi$ be as in the statement and, looking for a contradiction, assume that $\varphi$ is of class $C^1$. As $H$ is finitely generated, using a trick attributed to Muller \cite{Muller} and Tsuboi \cite{Tsuboi} (see also \cite{BMNR}), we~can assume that $\varphi(k)'(0)=\varphi(k)'(1)=1$ for every $k\in H$.

Let $\tau\colon(0, 1)\to \R$ be a non-decreasing map that semi-conjugates $\varphi$ to $\psi$, in~the sense that $\tau\varphi(k)=\psi(k)\tau$ for every $k\in H$. As Plante-like actions are minimal (Proposition~\ref{prop Plante-like minimal}), the semi-conjugacy $\tau$ must be continuous. To avoid confusion, we~will use Latin letters to denote points of $\R$ (where the Plante-like action $\eta$ is defined), and Greek letters for points of $[0, 1]$ (where the action $\varphi$ is defined). Set $J_n=\tau^{-1}(I_n)$ for $n\in \Z$. Note that $\varphi(\tilde{g})(J_n)=J_{n+1}$ for every $n\in \Z$. Let $\lambda:=\min_{\xi\in[0,1]}|D\varphi(\widetilde{g}^{-1})(\xi)|$. The mean value theorem implies that for every $n\in \Z$ we have the length estimate
\begin{equation} |J_{-n}|\geq |J_0|\lambda^{n}. \label{e-Jn}\end{equation}
We fix $N>\sfrac{1}{\lambda}$ and $\varepsilon>0$ such that
\begin{equation} (1-\varepsilon)^{2N+2}N\lambda>1. \label{e-epsilon-choice}
\end{equation}
Let $\sigma\in (0,1)$ be such that for any $\xi\in (\sigma, 1]$ and $s\in \{\widetilde{h}_0, \widetilde{g}, \widetilde{g}^{-1} \}$, one has $ D\varphi(s)(\xi)>1-\varepsilon$. Since the subgroup $L$ acts without fixed points in the Plante-like action, we~can find an element $a\in L$ such that $\eta(a)(\overline{I_0})\subset (\tau(\sigma), 1]$. Hence, choosing a preimage $\widetilde{a}\in \rho^{-1}(a)$ we have $\varphi(\widetilde{a})(\overline{J_0})\subset (\sigma, 1]$.

Fix $n\geq 1$. Given an $n$-tuple of integers $\underline{i}=(i_0, \ldots, i_{n-1})\in \{1,\ldots,N\}^n$, set $\widetilde{f}_{\underline{i}}=\widetilde{h}_{-n+1}^{i_{n-1}}\cdots \widetilde{h}_{-1}^{i_1} \widetilde{h}_0^{i_0}$ and ${f}_{\underline{i}}=\rho (\widetilde{f}_{\underline{i}})$. Note that the intervals of the form $\varphi(\widetilde{f}_{\underline{i}}\widetilde{a})(J_{-n})$ are pairwise disjoint when $\underline{i}$ varies. Indeed, the semi-conjugacy $\tau$ maps each such interval to the corresponding interval $\eta(f_{\underline{i}}a)(I_{-n})=\eta(af_{\underline{i}})(I_{-n})$, and these are pairwise disjoint by Lemma \ref{l-disjoint-intervals}, since the elements $f_{\underline{i}}$ all belong to the subgroup $L_{-n+1}$. We~shall estimate from below the size of $\varphi(\widetilde{f}_{\underline{i}}\widetilde{a})(J_{-n})$. For this, note that $\widetilde{f}_{\underline{i}}$ may be rewritten as
\begin{align*}
\widetilde{f}_{\underline{i}}&=\bigl (\widetilde{g}^{-n+1}\widetilde{h}_0^{i_{n-1}}\widetilde{g}^{n-1}\bigr)\cdots\bigl (\widetilde{g}^{-2}\widetilde{h}_0^{i_2}\widetilde{g}^2\bigr)\bigl (\widetilde{g}^{-1}\widetilde{h}_0^{i_1}\widetilde{g}\bigr)\widetilde{h}_0^{i_0}\\
&=\widetilde{g}^{-n+1}\widetilde{h}_0^{i_{-n+1}}\widetilde{g}\widetilde{h}_0^{i_{-n+2}}\widetilde{g}\cdots \widetilde{g}\widetilde{h}_0^{i_1}\widetilde{g}\widetilde{h}_0^{i_0}.
\end{align*}
Consider the sequence of intervals obtained by applying successively the terms in the latter expression to the interval $\varphi(\widetilde{a})(J_{-n})$. Observe that any such interval stays inside $(\sigma, 1]$: indeed, each application of the generator $\widetilde{g}$ or $\widetilde{h}_0$ moves the interval to the right, and the final term $\widetilde{g}^{-n+1}$ moves it to the left but is compensated by the previous $n-1$ instances of $\widetilde{g}$ with a positive power. By~the mean value theorem and the choice of $\sigma$, at each application of $\widetilde{g}, \widetilde{g}^{-1}$ or $\widetilde{h}_0$, the size of the interval may decrease by at most a factor $1-\varepsilon$. Hence
\[\bigl |\varphi(\widetilde{f}_{\underline{i}}\widetilde{a})(J_{-n})\bigr |\geq (1-\varepsilon)^{nN+2n-2}\left |\varphi(\widetilde{a})(J_{-n})\right |\geq C|J_0|(1-\varepsilon)^{nN+2n-2}\lambda^n,
\]
where we have used \eqref{e-Jn}, and the constant $C$ is $\min_{\xi\in[0,1]} D\varphi(\widetilde{a})(\xi)$. Summing over~$\underline{i}$ and using that the intervals are disjoint, we~have
\[1\geq \sum_{\underline{i}\in\{1,\ldots,N\}^n}\bigl |\varphi(\widetilde{f}_{\underline{i}}\widetilde{a})(J_{-n})\bigr |\geq C|J_0| (1-\varepsilon)^{-2} \bigl((1-\varepsilon)^{2N+2}N\lambda\bigr)^n,
\]
which is impossible, since by \eqref{e-epsilon-choice} the right-hand side is unbounded as $n$ tends to $+\infty$. This is the desired contradiction.
\end{proof}
\begin{thebibliography}{BGR23}
\bibitem[KKM19]{KKMj} S.-h. Kim, T. Koberda \& M. Mj – Flexibility of group actions on the circle, Lect. Notes in Math., vol. 2231, Springer, Cham, 2019.
\bibitem[BMBRT21]{BMRT} J. Brum, N. Matte Bon, C. Rivas \& M. Triestino – “Locally moving groups and laminar actions on the line”, 2021, to appear in Astérisque, arXiv:2104.14678.
\bibitem[Pla83]{Plante} J. F. Plante, “Solvable groups acting on the line”, Trans. Amer. Math. Soc. 278 (1983), no. 1, p. 401–414.
\bibitem[Mul82]{Muller} M.-P. Muller – “Sur l’approximation et l’instabilité des feuilletages”, Unpublished text, 1982.
\bibitem[Tsu84]{Tsuboi} T. Tsuboi – “Γ1 -structures avec une seule feuille”, in Transversal structure of foliations (Toulouse, 1982), Astérisque, no. 116, Société Mathématique de France, Paris, 1984, p. 222–234.
\bibitem[BMNR17]{BMNR} C. Bonatti, I. Monteverde, A. Navas \& C. Rivas – “Rigidity for C 1 actions on the interval arising from hyperbolicity I: solvable groups”, Math. Z. 286 (2017), no. 3-4, p. 919–949.
\bibitem[CR16]{ClayRolfsen} A. Clay \& D. Rolfsen – Ordered groups and topology, Graduate Studies in Math., vol. 176, American Mathematical Society, Providence, RI, 2016.
\bibitem[Nav11]{Navas-book} A. Navas, Groups of circle diffeomorphisms, Chicago Lectures in Math., University of Chicago Press, Chicago, IL, 2011.
\end{thebibliography}
\end{document}
